\section{Introduction}\label{sec:intro}
\owner{Wei Yew (draft); Ky to review}

\subsection{Motivation}
Clustering is used where no answer key exists. A seismologist grouping earthquake epicentres wants
active fault systems; a housing analyst grouping resale transactions wants market segments. Neither
has labels to check a result against.

Each clustering method carries its own idea of what a cluster is. K-means~\cite{lloyd1982} looks for
compact groups around a centre and assigns every point to one of them. DBSCAN~\cite{ester1996} and
HDBSCAN~\cite{campello2013} look for dense regions separated by sparse space, and leave points in the
sparse space as noise. Whether or not that idea matches the data, the method returns a partition. On
one year of global earthquakes, K-means returns six tidy regions, one of which joins five plate
boundaries lying 2{,}932 to 11{,}449\,km apart (Section~\ref{sec:shape}). Nothing in the output says
that this has happened. An analyst without labels usually judges the result by an internal score such
as the silhouette~\cite{rousseeuw1987}, which rates how compact and well separated the clusters are,
and that score is itself built on the K-means idea of a cluster.

\subsection{Research question}
This report asks one question: \emph{when a clustering method's idea of a cluster does not match the
data, does it fail, and can the failure be detected without labels?} We split it into three parts.
\begin{description}\itemsep1pt
\item[RQ1, cause.] Which property of the data breaks which method? We consider cluster shape, density
that varies between clusters, the number of dimensions, the absence of gaps between groups, feature
scale, and values recorded in steps.
\item[RQ2, detection.] Do checks that need no labels reveal those failures? We test the silhouette, the
Davies--Bouldin index, the share of points labelled noise, and a fake-data check: the same method run
on data whose features have been shuffled independently, which keeps each feature's values but removes
any joint structure.
\item[RQ3, remedy.] For each failure, does a known remedy help, and what does it cost?
\end{description}

We compare four methods, K-means with random starts, K-means++~\cite{arthur2007}, DBSCAN and HDBSCAN,
on eight datasets (Table~\ref{tab:data}). Each dataset was chosen to test one data property, most of
them a property that breaks at least one method's assumption. Six have class labels, so a method's output can be scored against them with the
Adjusted Rand Index (ARI)~\cite{hubert1985}, and the internal scores can be checked against that
answer. Two have no labels (earthquakes, Uber pickups); there we judge failures on maps, against known
geography.
Gaussian mixtures, spectral clustering and agglomerative clustering appear only as reference points on
the labelled datasets where we ran them (D1, D5, D6, D7), not as part of the main comparison.

\subsection{Findings in brief}
% Sources: lam/results/{make_moons,fashion_mnist,wine}_results.csv; hoanganh/results/breast_cancer_results.csv;
% SC4020-Ky/data/raw/main.csv (Ecoli); weiyew/sc4020-earthquake/results/{case_density_regions,
% method_comparison,ablation_distance}.csv; weiyew/sc4020-hdb/results/comparison.csv.
% RQ2 count, method with top silhouette vs method with top ARI, density settings within the 50% noise cap:
% D1 K-means 0.495 (ARI 0.481) vs Spectral 1.000; D4 HDBSCAN 0.483 (0.339) vs K-Means++ 0.731;
% D5 K-means 0.202 (0.343) vs Spectral 0.398; D6 HDBSCAN 0.359 (0.263) vs K-means 0.897;
% D7 K-means 0.343 (0.654) vs GMM 0.774; D8 K-means is top on both (not counted).
\begin{description}\itemsep1pt
\item[RQ1.] Each method failed where the data broke its assumption. K-means split the two interleaved
half-moons of D1 wrongly (ARI 0.48) and merged distant belts on D2. With one radius, DBSCAN clustered
90\% of earthquakes in the dense south-west Pacific but only 20\% in the sparser Americas. In 50
dimensions (D5) it barely matched the labels (ARI 0.03). On D8, where housing sales form one mass with no
gaps, both density methods returned one dominant cluster.
\item[RQ2.] The silhouette did not reveal these failures. On five of the six labelled datasets, the
method with the highest silhouette was not the one that matched the labels best. On D1 it ranks K-means
first (0.495 against 0.392 for DBSCAN), although DBSCAN recovers the moons (ARI 0.98 against 0.48). On
Ecoli (D4) it orders the four methods almost exactly opposite to ARI (Spearman $\rho=-0.80$). The
fake-data check, which we ran on D2 and D8, does better, but only partly. It detects when a method finds nothing real: on D8 both
density methods score below the same setting on shuffled data. It does not detect the wrong structure:
on D2, K-means scores well above shuffled coordinates (silhouette 0.528 against 0.391) while merging
unrelated belts. Two failures were found only by checks aimed at a known problem: clusters cut at the
$180^\circ$ meridian, and clusters that are just one recorded storey band.
\item[RQ3.] Known remedies repaired some failures and not others. Measuring distance on the sphere kept
all 92 close earthquake pairs across the $180^\circ$ meridian together, where flat coordinates kept
none. HDBSCAN raised the clustered share of the Americas from 20\% to 62\%, at the cost of trimming the
densest belts (78\% against 90\%). Nothing we tried produced meaningful density clusters on D8 with all
four features: capping the size of the largest cluster only made the rule select storey bands.
\end{description}
\todo{Lam, Ky, Hoang Anh}{confirm the RQ2 numbers above are your settings under the shared selection
rule (Section~\ref{sec:protocol}); the comment above lists each one.}

Section~\ref{sec:methods} describes the methods and Section~\ref{sec:experiments} the datasets,
protocol and results, grouped by data property. Section~\ref{sec:conclusion} turns them into guidance.

\subsection{Related work}
\paragraph{Centroid and mixture models.} Lloyd's algorithm~\cite{lloyd1982}
minimises the within-cluster sum of squared distances by alternating assignment and
centroid updates. It converges to a local optimum that depends on the starting
centres. K-means++~\cite{arthur2007} picks each starting centre with probability
proportional to its squared distance from the centres already chosen, which gives an
$O(\log k)$ approximation to the optimum in expectation. A Gaussian mixture fitted by
EM~\cite{dempster1977} generalises K-means to ellipsoidal clusters with soft
assignments. Agglomerative clustering builds a hierarchy by repeatedly merging the
closest pair of clusters. Under Ward's criterion~\cite{ward1963} the merge cost is the
increase in within-cluster variance, the same objective as K-means.

\paragraph{Density-based methods.} DBSCAN~\cite{ester1996} defines a cluster as a
maximal set of density-connected points. A point is a core point if at least
\texttt{minPts} points lie within radius $\varepsilon$. Points reachable from no core
point are noise. The method needs no $k$ and imposes no shape, but its single global
$\varepsilon$ assumes all clusters have similar density. OPTICS~\cite{ankerst1999}
orders points by reachability so that clusters at several densities can be read from
one plot. HDBSCAN~\cite{campello2013} builds a hierarchy over mutual reachability
distances and keeps the clusters that persist longest. This removes $\varepsilon$ and
leaves a minimum cluster size as the main parameter.

\paragraph{Graph-based methods.} Spectral clustering~\cite{ng2002} embeds the points
using the leading eigenvectors of a normalised affinity matrix, then runs K-means in
that embedding. Groups connected through the affinity graph become compact in the
embedding, so non-convex clusters can be separated. A dense affinity matrix needs
$O(n^2)$ memory, which limits the method to modest $n$.

\paragraph{Validation and cluster tendency.} Internal indices, such as the
silhouette~\cite{rousseeuw1987}, Davies--Bouldin~\cite{davies1979} and
Calinski--Harabasz~\cite{calinski1974} indices, score a partition from the data
alone. External indices such as the Adjusted Rand Index~\cite{hubert1985} compare it
with known labels. Von Luxburg et al.~\cite{luxburg2012} argue that no criterion can
judge a clustering apart from its intended use. A separate line of work tests whether
data contain clusters at all. The Hopkins statistic~\cite{hopkins1954} compares
nearest-neighbour distances in the data with those in uniformly random points. The
gap statistic~\cite{tibshirani2001} compares within-cluster dispersion with its value
on reference data drawn from a null distribution. Our fake-data check applies the
same idea to every method we compare. In high dimensions the ratio of nearest to
farthest neighbour distance tends to one~\cite{beyer1999}, which weakens any method
built on a density threshold.
